\chapter*{Résumé}
\addcontentsline{toc}{chapter}{Résumé}

Dans le cadre de mon stage de licence de deux mois au sein de 
l'entreprise Bewize, j'ai développé une fonctionnalité complète 
de recueil de feedbacks et de réclamations pour l'application 
mobile Bewize. Ce projet avait pour objectif de permettre aux 
utilisateurs de l'application mobile de transmettre facilement 
leurs retours directement depuis l'application, et à l'équipe 
interne Bewize de les consulter et les traiter via un dashboard 
monitor dédié.

La solution développée repose sur une architecture en quatre 
composants principaux : une page WebView développée en React JS 
contenant le formulaire de saisie, intégrée dans l'application 
mobile React Native via le composant WebView, un back-end développé 
avec Spring Boot exposant des APIs REST pour la gestion des données, 
et un dashboard monitor en React JS permettant à l'équipe de 
consulter et traiter les retours reçus. Les données sont persistées 
dans une base de données PostgreSQL et l'ensemble des services 
est conteneurisé avec Docker.

Le projet a été réalisé en suivant la méthodologie Agile Scrum, 
organisé en trois Sprints de deux semaines chacun, sous la 
supervision du Product Manager et avec l'encadrement technique 
du Tech Lead de Bewize. Les tests réalisés à trois niveaux 
(unitaires, intégration et fonctionnels bout en bout) ont confirmé 
la qualité et la fiabilité de la solution, avec un taux de réussite 
de 100\% sur l'ensemble des 39 tests effectués.

Ce stage m'a permis d'acquérir des compétences techniques solides 
en développement web et mobile, de découvrir le travail en équipe 
dans un environnement agile professionnel, et de comprendre le 
lien essentiel entre le développement produit et les besoins 
métier d'une entreprise numérique.

\vspace{1cm}

\textbf{Mots-clés :} React JS, React Native, Spring Boot, PostgreSQL, 
Docker, Agile Scrum, WebView, API REST, Feedback, Réclamation, 
Dashboard monitor, Bewize.